
Der Ressourcen-Fit beschreibt die Passung zwischen dem Ressourcenbedarf einer
ERP-Einführung und den verfügbaren finanziellen, personellen, zeitlichen und
fachlichen Kapazitäten eines KMU. ERP-Systeme können zwar langfristig dazu
beitragen, vorhandene Ressourcen effizienter zu nutzen, ihre Einführung bindet
jedoch zunächst selbst erhebliche Ressourcen. Damit wird Ressourcen-Fit zu
einer zentralen Voraussetzung dafür, ob ein KMU ein ERP-Projekt realistisch
planen, durchführen und anschließend betreiben kann.

Finanzielle Ressourcen betreffen dabei nicht nur die unmittelbaren
Anschaffungskosten eines ERP-Systems. Neben Softwarelizenzen, Hardware oder
Cloud-Abonnements entstehen häufig weitere Kosten für externe Beratung,
Systemanpassungen, Datenmigration, Schulungen, Tests und laufenden Support
\parencite{Jha2008, Leyh2015}. Gerade diese indirekten oder versteckten Kosten können
für KMU kritisch sein, da sie häufig über geringere finanzielle Puffer
verfügen als Großunternehmen \parencite{Leyh2015}. Ein niedriger finanzieller Ressourcen-Fit liegt
daher vor, wenn ein ERP-Projekt nur anhand der reinen Systemkosten bewertet
wird und zusätzliche Aufwände für Einführung, Anpassung und Betrieb nicht
realistisch eingeplant werden.

Auch personelle Ressourcen sind für ERP-Projekte entscheidend. Die Einführung
erfordert ein internes Projektteam, das fachliches Wissen aus den relevanten
Unternehmensbereichen einbringt und als Schnittstelle zwischen Unternehmen,
Anwendern und externen Beratern fungiert \parencite{Leyh2015}. Besonders wichtig sind dabei
Mitarbeitende, die die bestehenden Prozesse gut kennen und Anforderungen an
das System formulieren können \parencite{Leyh2017}. Für KMU ist dies problematisch, weil gerade
diese Wissensträger häufig zugleich stark in das operative Tagesgeschäft
eingebunden sind. Werden sie nicht ausreichend für Projektaufgaben entlastet,
kann dies die Qualität von Prozessanalyse, Tests, Schulung und Einführung
beeinträchtigen \parencite{Jha2008,Leyh2015}.

Zeitliche Ressourcen bilden eine weitere Dimension des Ressourcen-Fits.
ERP-Einführungen umfassen mehrere zeitintensive Phasen, darunter
Ist-Analyse, Softwareauswahl, Prozessanalyse, Datenmigration,
Systemkonfiguration, Testphase, Schulung, Go-Live und Nachbetreuung \parencite{Leyh2015,Leyh2019}. Diese
Aktivitäten laufen in KMU häufig parallel zum Tagesgeschäft. Dadurch entstehen
Opportunitätskosten, da Zeit für Projektarbeit nicht gleichzeitig für
operative Tätigkeiten genutzt werden kann \parencite{Leyh2015,Leyh2017}. Ein hoher Ressourcen-Fit setzt
daher voraus, dass für zentrale Projektaufgaben ausreichend Zeit eingeplant
wird und Belastungsspitzen während der Einführung berücksichtigt werden.

Neben Geld, Personal und Zeit ist auch fachliches Wissen erforderlich.
ERP-Projekte setzen IT-Know-how, Prozessverständnis,
Projektmanagementkompetenz und Kenntnisse über Datenmigration und
Systemanpassung voraus. Gerade KMU verfügen jedoch häufig nur über begrenzte
interne IT-Kapazitäten und weniger Erfahrung mit komplexen
Digitalisierungsprojekten \parencite{Leyh2015,Jha2008}. Dadurch steigt die
Abhängigkeit von externen Beratern oder Softwareanbietern. Diese Abhängigkeit
ist besonders dann kritisch, wenn das Unternehmen seine eigenen Anforderungen
nicht ausreichend präzise beschreiben oder externe Vorschläge nicht
angemessen bewerten kann \parencite{Leyh2015,Leyh2017}.

Ein hoher Ressourcen-Fit liegt somit vor, wenn ein KMU die ERP-Einführung und
die spätere Nutzung finanziell, personell, zeitlich und fachlich tragen kann.
Dazu gehören ein realistisches Budget, verfügbare interne Wissensträger,
ausreichende Zeit für Projektarbeit, Tests und Schulungen sowie Zugang zu
notwendiger technischer und methodischer Expertise. Für KMU ist diese
Dimension besonders kritisch, weil fehlende Ressourcen nicht nur einzelne
Projektaktivitäten erschweren, sondern den gesamten Nutzen eines ERP-Systems
gefährden können.

